% !TEX root = ../macro.tex
\section{Growth Models: Mechanisms, Comparisons, and Reality}\label{sec:growth-synthesis}

A region builds factories, trains workers, starts research laboratories, and
eventually faces slower growth and an older population. Which growth model
describes it? Several may help, because each isolates a different question.
The useful comparison is between mechanisms, rather than between labels for
countries or stages of development.

This chapter follows that region through a sequence of questions. The region
is a teaching construction, not a disguised account of a particular country.
We add mechanisms as the questions require them, keep their assumptions
visible, and use historical evidence at specific points to connect the
construction to actual experience. The aim is to learn how to move from an
observation to a model and then back to an observation that could test it.

\subsection{A map of the questions, not a ranking of the models}

Growth accounting asks where measured output growth appears: capital,
labor, or a residual. A growth model asks why those quantities change.
Harrod and Domar ask whether spending and capacity expansion remain
compatible. Solow asks how capital intensity adjusts when inputs can
substitute. Ramsey asks who gives up consumption to finance accumulation.
OLG changes who is alive, who owns assets, and which generations can trade.
Endogenous-growth models ask what keeps the return to reproducible inputs
or new knowledge from fading.

\Needspace{23\baselineskip}
\noindent\textbf{Production, saving, and ownership.}
\begin{center}
\small
\begin{tabular}{p{0.17\linewidth}p{0.32\linewidth}p{0.37\linewidth}}
\hline
Framework & Question it closes & Situation that makes the question concrete\\
\hline
Growth accounting & Which measured inputs and residual accompany growth? &
Output rises, but employment and capital also rise. How much is left to explain?\\[4pt]
Harrod--Domar & Which investment path keeps demand, capacity, and plans compatible? &
Building factories creates spending today and capacity that needs customers later.\\[4pt]
Solow & How does capital per effective worker adjust under a given saving rule? &
More machines initially help a capital-scarce region; their marginal contribution then falls.\\[4pt]
Ramsey & How do households choose consumption and saving over time? &
Higher expected returns make postponing consumption more attractive, but the sacrifice has a cost.\\[4pt]
OLG & How do life-cycle saving and trade between living cohorts determine assets? &
Workers buy retirement assets from older owners; future workers cannot sign contracts today.\\
\hline
\end{tabular}
\end{center}

\Needspace{23\baselineskip}
\noindent\textbf{Skills, ideas, and the use of resources.}
\begin{center}
\small
\begin{tabular}{p{0.17\linewidth}p{0.32\linewidth}p{0.37\linewidth}}
\hline
Mechanism & What must be explained & Situation that makes the question concrete\\
\hline
AK & Why can a reproducible stock retain its marginal return? &
Machines, skills, or other complementary inputs may expand together; the aggregate linearity still needs a justification.\\[4pt]
Human capital & How much time goes to work rather than learning? &
A technician leaves production for training, reducing output now to improve future capability.\\[4pt]
Innovation & Who pays to create a reusable idea, and who receives its return? &
A firm funds a design team; rivals may benefit from knowledge that the firm cannot fully charge for.\\[4pt]
Semi-endogenous growth & How does research effort interact with the difficulty of discovery? &
More researchers can be needed just to maintain a previous pace of technical improvement.\\[4pt]
Allocation & Which producers use the existing inputs? &
A productive workshop cannot expand while a less productive firm retains equipment or financing.\\
\hline
\end{tabular}
\end{center}

These are distinct dimensions. OLG does not follow Ramsey because it has
more technology; it changes the household structure. An OLG economy can
also innovate. A Ramsey economy can face demand constraints once the
required market mechanisms are added. An allocation problem can matter
at low or high income. Combining them requires a new, consistent equilibrium,
not the addition of conclusions obtained under incompatible assumptions.

\Needspace{11\baselineskip}
\begin{greenblock}
\textbf{Add the mechanism that answers the next question.}
Start with what is scarce, who chooses its use, and what prevents adjustment.
Then ask whether the unresolved issue concerns demand, saving, generations,
skills, discovery, or allocation. A later model in these notes is not a claim
that every economy must pass through the same development sequence.
\end{greenblock}

\subsection{A region builds factories: demand, capacity, and diminishing returns}

\subsubsection{First thought: more investment should mean more growth}

Suppose the region has workers, basic skills, and access to an existing
production method, but few machines. A factory project employs construction
workers and buys equipment. Once completed, it can produce goods for years.
It is natural to start with investment as the driver of growth. Two questions
are already hidden in that statement: will the new capacity be used, and how
much output will another unit of equipment add?

Harrod--Domar makes the first question visible. In Domar's experiment,
\[
\dot Q=\kappa I,\qquad Y^d=I/s,\qquad \dot Y^d=\dot I/s,
\]
where $Q$ is capacity, $Y^d$ is demand-supported output, $s$ is a fixed
saving propensity, and $\kappa$ is a constant capacity increment per unit of
investment. With no depreciation in this experiment, constant positive
investment keeps adding capacity but leaves demand-supported output
unchanged. Maintaining an initially matched path requires
$\dot I/I=s\kappa$. Section~\ref{sec:harrod-domar} derives the conditions;
the relation is a compatibility requirement, not a forecast of investment.

Harrod adds a related planning question: does realized growth justify the
capacity that firms intended to install? With fixed saving propensity $s$
and desired capital--output ratio $v_r>0$, his no-depreciation benchmark
gives warranted growth $g_w=s/v_r$. Labor and labor productivity instead
set natural growth $g_n=n+g_A$. Matching these rates is an additional
condition, not an identity. If sales disappoint and firms cut investment,
the resulting loss of spending can further undermine sales. That feedback
depends on the stated investment response; fixed proportions alone do not
prove instability.

The real counterpart is a new industrial park whose construction creates
business for local suppliers, but whose finished factories still need buyers.
A factory is a productive asset; construction spending is a flow of demand.
Confusing them makes it easy to count the spending effect and silently assume
that all future output will be sold. Whether customers are domestic or foreign,
and whether prices or credit adjust, requires additional mechanisms.

\subsubsection{Second thought: can every new factory be as productive as the first?}

Now suppose demand is sufficient and markets clear. The region can alter
its mix of workers and equipment. Solow isolates this different question:
with $k=K/(AL)$ and $y=Y/(AL)$,
\[
y=f(k),\qquad \dot k=sf(k)-(n+g_A+\delta)k.
\]
Here $A$ is labor efficiency, $L$ grows at rate $n$, $A$ grows at an
exogenous rate $g_A$, and $\delta$ is depreciation. Diminishing $f'(k)$ means
that another machine adds less output when each effective worker already
has more equipment. Capital intensity adjusts rather than requiring one
fixed capital--output ratio to remain compatible with every growth path.

For a numerical check, hold $A$ and $L$ fixed at a date and let
$f(k)=k^{1/3}$. A 20\% increase in $k$ raises $y$ by
\[
\frac{y_1}{y_0}=1.2^{1/3}\simeq1.0627,
\qquad
\frac{f'(k_1)}{f'(k_0)}=1.2^{-2/3}\simeq0.8855.
\]
Output rises by about 6.27\%, while the marginal product falls by about
11.45\%. There is no contradiction: total output and the return to the next
unit of capital are different objects. These are illustrative quantities,
not estimates of an industrialization program.

Solow does not prove that newly built capacity always finds buyers.
Harrod--Domar does not prove that substitution and price adjustment never
occur. The region can have a short-run utilization problem and a long-run
capital-intensity problem at the same time. The models differ in the
adjustment mechanism they choose to expose.

\subsection{Someone must save: from an investment target to household choices}

\subsubsection{A higher eventual income is not a free transition}

Suppose officials propose raising the saving rate from $0.20$ to $0.30$.
In the continuous-time Cobb--Douglas Solow benchmark,
\[
k^*=\left(\frac{s}{n+g_A+\delta}\right)^{1/(1-\alpha)},
\qquad
y^*=\left(\frac{s}{n+g_A+\delta}\right)^{\alpha/(1-\alpha)}.
\]
With $\alpha=1/3$, unchanged technology growth, and
$n+g_A+\delta>0$, the eventual ratio of output per effective worker is
\[
\frac{y^*_{\mathrm{new}}}{y^*_{\mathrm{old}}}
=\left(\frac{0.30}{0.20}\right)^{1/2}\simeq1.2247.
\]
But capital is inherited at the policy date. If the saving change is imposed
immediately, output initially stays fixed and consumption changes by
\[
\frac{c(0^+)_{\mathrm{new}}}{c(0^-)_{\mathrm{old}}}
=\frac{1-0.30}{1-0.20}=0.875.
\]
An eventual output gain of about 22.47\% comes with an initial consumption
cut of 12.5\%. This calculation compares two imposed saving rules. It does
not establish the welfare ranking of their entire consumption paths.

The next question is unavoidable: whose consumption falls, and why would
they agree? Ramsey replaces the saving instruction with preferences,
resources, prices, and a choice. In the fixed-population, fixed-technology,
competitive CRRA benchmark, with $c$ and $k$ per person,
\[
\frac{\dot c}{c}=\frac{f'(k)-\delta-\rho}{\theta},
\qquad \rho>0,\quad\theta>0.
\]
A unit not consumed today can earn a net return. Consumption grows when
that return exceeds impatience, with $\theta$ governing the response.
The Euler condition must still be combined with the resource constraint,
initial capital, feasibility, and the terminal condition. It is not a saving
rate by itself.

In both the Solow and standard Ramsey benchmarks, continued growth of
output per worker on a balanced path still comes from the specified growth
of labor efficiency. Optimizing saving changes the accumulation path and
its welfare interpretation; it does not remove that technological assumption.
The OLG benchmark changes generations rather than supplying a new technology
engine. This is the gap addressed by endogenous growth.

For the region, a credible improvement in future production opportunities
may change desired saving. A temporary income loss may instead encourage
households to use assets to protect consumption. Taxes, borrowing limits,
and uncertainty can alter both responses. The basic Ramsey model supplies
the intertemporal trade-off; it does not establish every household's response
when its omitted constraints become important.

\subsubsection{The same residents are also different generations}

Now let the region have young workers, retirees, and future entrants.
The question is no longer only how one continuing household spreads
consumption over time. Workers must acquire claims that will pay in
retirement, while current retirees sell assets or receive transfers.

In the two-period OLG benchmark, cohort $t$ chooses $c_{1,t}$ and
$c_{2,t+1}$ with lifetime utility $u(c_{1,t})+\beta u(c_{2,t+1})$,
where $\beta>0$. Without taxes, transfers, or bequests, its budgets are
\[
c_{1,t}+a_{t+1}=w_t,
\qquad c_{2,t+1}=R_{t+1}a_{t+1}.
\]
Here $a_{t+1}$ is saving per young worker, $w_t$ is labor income, and
$R_{t+1}$ is the gross payoff on saving. A worker's saving can become the
next period's physical capital, but the economy's net saving also reflects
retirees' decumulation. The number and resources of each cohort matter.

A pension reform therefore has several margins: contributions reduce
workers' disposable income, expected benefits may replace private retirement
saving, and the first recipients and transition cohorts face different
budgets. Fewer future workers can also lower a pay-as-you-go system's
support ratio. None of these observations alone fixes the effect of aging
on aggregate saving or capital: longevity, retirement ages, wages, transfers,
and behavior must be specified.

\Needspace{11\baselineskip}
\begin{blueblock}
\textbf{An aging economy is not automatically dynamically inefficient.}
In the deterministic, stationary OLG benchmark with no technology growth,
overaccumulation means $r<n$, equivalently $R<1+n$, under the chapter's
depreciation convention. Age structure or a low observed safe interest rate
alone does not verify that condition. Risky capital returns, government
claims, and a changing population are not interchangeable objects.
\end{blueblock}

Ramsey and OLG thus offer different ownership and welfare comparisons for
the same investment program. Ramsey tracks a continuing household's
consumption path. OLG asks which living and future cohorts gain or lose,
including the initially old. Greater eventual output need not compensate
those who finance the transition.

\subsection{The next machine needs a skill; the next skill needs an idea}

\subsubsection{A machine shortage becomes a capability question}

The region now has equipment, but technicians cannot maintain it, managers
cannot coordinate production, and workers need time to master new methods.
Buying the same machine again may not solve these problems. What looked
like a capital shortage can include missing complementary capabilities.

Human capital makes one part explicit. In the two-sector benchmark,
\[
Y=K^\alpha(vH)^{1-\alpha},\qquad
\dot H=[\eta(1-v)-\delta_H]H,
\]
where $H$ is the skill stock, $v\in[0,1]$ is the fraction of time spent
working, $1-v$ is learning time, $0<\alpha<1$, $\eta>0$ is learning
productivity, and $\delta_H\geq0$ is skill depreciation. Moving time into
learning reduces current production but raises future productive capability.
The model's two shadow values put that opportunity cost next to its benefit.

This helps distinguish three interventions. More training hours changes
$1-v$; more effective instruction changes $\eta$; a better inherited skill
base changes $H(0)$. They are different experiments. School enrollment or a
training budget does not directly measure acquired skills or learning
productivity. Credit constraints, teaching quality, and employer incentives
would need to be added to explain those gaps.

AK provides a useful warning at this point. Writing $Y=BK$ assumes that
the reproducible aggregate retains its marginal product. It does not explain
why installing more identical machines is sufficient. Joint skill and capital
accumulation can provide one mechanism, under the conditions derived in
Section~\ref{sec:endogenous-growth}. An aggregate externality provides
another, but then private and social returns differ. The equation alone
cannot choose between them.

\subsubsection{Learning an existing method differs from creating a new one}

Suppose local engineers learn a production method already used abroad.
This improves local productivity without expanding the world's stock of
ideas. Adoption requires capability and often costly adaptation; the
distance from existing best practice matters. At the frontier, the engineers
must instead discover a method that is not yet available.

Korea provides a historical connection. A 2007 World Bank--Korea Development
Institute study describes learning through imported equipment, production
arrangements, training, and reverse engineering, followed by a greater role
for domestic innovation.\footnote{Joonghae Suh and Derek H. C. Chen, eds.
(2007), \textit{Korea as a Knowledge Economy: Evolutionary Process and
Lessons Learned}, World Bank, especially Chapters 6--8 and pp.~168--169
on technological adoption and innovation.
\href{https://documents1.worldbank.org/curated/en/923361468327364601/pdf/409300PAPER0KR101OFFICIAL0USE0ONLY1.pdf}{World Bank report}.}
This illustrates why equipment, absorptive skills, and research can be
complements. It does not identify one model as the complete explanation of
Korean growth, or turn the sequence into a causal estimate. Trade, finance,
institutions, and policy also belong in that account.

Return to our region. A design team uses workers who could have produced
goods today. Once created, an idea may be usable by many producers without
being depleted. A patent or another means of earning revenue can fund the
team, but its private return need not capture all future knowledge benefits.
This is the additional question in the innovation model: research effort
must be supported by expected rewards and market entry, not just entered
as a positive term in a production function.

\Needspace{11\baselineskip}
\begin{greenblock}
\textbf{Three objects that a productivity residual can hide.}
Adoption brings an existing method into local use. Human capital enables
people to use or develop methods. Innovation adds new ideas. A measured
productivity improvement may contain all three, as well as allocation and
measurement effects. Growth accounting locates a residual; it does not
separate these mechanisms without further information.
\end{greenblock}

\subsubsection{More researchers need not mean permanently faster growth}

In the earlier linear research technology, $\dot A/A=\eta L_R$, a
permanently larger research workforce $L_R$ permanently raises idea growth.
Changing the technology to
\[
\dot A=\eta A^\phi L_R^\lambda,
\qquad \phi<1,\quad0<\lambda\leq1,
\]
changes that conclusion. With a fixed positive research share and labor
growing at $n>0$, a balanced path instead has
$g_A=\lambda n/(1-\phi)$. A larger research share then raises the level
path and changes the transition, while this balanced growth rate is unchanged.
The calculation is conditional on the research share; it is not a new
solution for optimal research entry.

Bloom, Jones, Van Reenen, and Webb (2020) examine research effort and
research productivity in historical aggregate and sectoral evidence. Their
finding that rising effort can accompany declining research productivity
motivates checking the research technology rather than assuming a constant
return per researcher.\footnote{Nicholas Bloom, Charles I. Jones, John
Van Reenen, and Michael Webb (2020), ``Are Ideas Getting Harder to Find?''
\textit{American Economic Review}, 110(4), 1104--1144.
\url{https://doi.org/10.1257/aer.20180338}.}
It does not prove that one value of $\phi$ explains every industry or period.
The practical question is how much effective research is needed to sustain
the next improvement, not simply how many laboratories are built.

\subsection{The same resources can produce more, but who gains?}

Even after new factories, training, and research, the region can waste
productive opportunities. Imagine two firms with different technologies.
One has idle equipment; the other could use additional equipment productively
but faces an explicit tax or financing obstacle. More aggregate investment
is only one possible response. Moving the existing stock can also raise
output, provided the gain exceeds movement, adjustment, and other costs.

The static allocation experiment in Section~\ref{sec:endogenous-growth}
holds total capital and complementary inputs fixed. With
$Y=2\sqrt{k_1}+\sqrt{k_2}$ and $k_1+k_2=1$, equal allocation gives
$Y=3/\sqrt{2}$; the efficient allocation $(0.8,0.2)$ gives $Y=\sqrt{5}$.
No new equipment or idea is required. The mechanism is a change in where
capital operates. Repeating an already completed reallocation does not
generate another permanent increase in the growth rate.

Hsieh and Klenow's study of manufacturing in China and India makes the
allocation question empirical by constructing counterfactuals under a
specified firm-level model.\footnote{Chang-Tai Hsieh and Peter J. Klenow
(2009), ``Misallocation and Manufacturing TFP in China and India,''
\textit{Quarterly Journal of Economics}, 124(4), 1403--1448.
\url{https://doi.org/10.1162/qjec.2009.124.4.1403}.}
The connection is the distinction between available technologies and the
allocation of inputs across producers. The counterfactual remains conditional
on how productivity, prices, markups, and wedges are measured. A difference
in observed revenue per unit of capital is not sufficient to diagnose a
credit distortion.

For residents, a productive reallocation can still create losers. Workers
may need to move, existing asset values can fall, and a firm owner may lose
a protected position. Our static calculation has neither relocation costs
nor a distributional account. Those costs must be introduced before claiming
that every observed obstacle is removable at no sacrifice.

\Needspace{11\baselineskip}
\begin{blueblock}
\textbf{Three comparisons that must stay separate.}
More output from given inputs is a production comparison. More consumption
after replacement and accumulation is a resource comparison. A preferred
consumption path, or an improvement for every cohort, is a welfare comparison.
Solow's Golden Rule, Ramsey welfare, and OLG Pareto improvements answer
different questions. None follows from an investment or GDP target alone.
\end{blueblock}

At the research frontier, allocation and growth can interact. Moving skilled
workers toward research may alter future knowledge creation, while raising
the rewards to incumbent firms may weaken entry. To sign a growth effect,
specify those margins and solve the resulting decisions. The static output
gain itself does not settle the dynamic incentive question.

\subsection{A growth slowdown: work backward from observations}

Suppose the region's GDP growth falls after years of expansion. One account
says it has exhausted capital accumulation; another blames aging; another
calls for more research. All are hypotheses. Start by asking what actually
slowed and which observations would distinguish the mechanisms.

\subsubsection{Separate production per person from the number of people working}

For total hours $L$, employment $E$, and population $N$,
\[
\frac{Y}{N}=\frac{Y}{L}\frac{L}{E}\frac{E}{N},\qquad
g_{Y/N}=g_{Y/L}+g_{L/E}+g_{E/N}.
\]
The identity does not require a growth model. Falling employment per resident
can reduce output per person even if output per hour continues improving.
Aging can affect that ratio, but participation and retirement behavior can
also change. A decline in total GDP growth could additionally reflect slower
population growth. Specify the denominator before choosing the explanation.

\Needspace{27\baselineskip}
\noindent\textbf{Different diagnoses of the same headline.}
\begin{center}
\small
\begin{tabular}{p{0.27\linewidth}p{0.28\linewidth}p{0.31\linewidth}}
\hline
Observation to investigate & Candidate mechanism & What would help distinguish it?\\
\hline
Capital per worker keeps rising; its additional output contribution falls &
Solow transition toward a conditional balanced path &
Capital quality, effective labor, returns, and the distance from the conditional target\\[4pt]
Output per person slows more than output per hour &
Employment, hours, and demographic composition &
Participation, age-specific employment, hours, and population measures\\[4pt]
Equipment is installed, but cannot be used effectively &
Skills, complementary inputs, or adoption costs &
Training outcomes, vacancies, maintenance, and technology use\\[4pt]
Research effort rises without a matching pace of improvement &
Research productivity, duplication, or measurement &
Effective research inputs and comparable measures of useful innovation\\[4pt]
Productive firms remain small while others retain inputs &
Allocation wedges, adjustment costs, or market power &
Firm prices, quantities, costs, risk, and a specified counterfactual\\[4pt]
Orders and utilization fall before installed capacity changes much &
Spending, expectations, or short-run adjustment &
Orders, inventories, hours, credit conditions, and the timing of capacity changes\\
\hline
\end{tabular}
\end{center}

The last column contains useful evidence, not automatic causal tests. For
example, low measured returns can reflect expectations or risk, and unfilled
vacancies can reflect wages as well as skills. A diagnosis becomes stronger
when a proposed mechanism explains several observations and survives a
counterfactual comparison.

\subsubsection{Separate using capacity from expanding it}

For a simple diagnostic illustration, define $Q$ as output capacity under
a stated operating convention and $u=Y/Q>0$ as utilization. Then
\[
Y=uQ,\qquad \Delta\log Y=\Delta\log u+\Delta\log Q.
\]
This is an identity once $Q$ is defined, not an independent measure of
capacity or a behavioral theory. If capacity grows by 3\% over a year but
utilization falls from $0.90$ to $0.81$, actual output changes by
\[
\frac{Y_1}{Y_0}=\frac{0.81}{0.90}\times1.03=0.927.
\]
Output falls by 7.3\% although capacity rises. If utilization subsequently
recovers, fast observed growth need not mean that the growth rate of capacity
has permanently increased. The chosen numbers are an accounting example.

Figure~\ref{fig:growth-path-diagnosis} compares three distinct changes in
output paths. They are deliberately constructed illustrations, not estimated
responses. In an actual economy, capacity and utilization are difficult to
separate and can affect one another: prolonged weak sales can discourage
investment, training, or entry. The short run can leave a long-run trace.

\begin{figure}[!ht]
\centering
\includegraphics[width=\linewidth]{growth_path_diagnosis.pdf}
\caption{A temporary utilization decline, a lasting level loss, and slower
trend growth produce different paths relative to the same benchmark. The
shock occurs at model date 10. Both panels show constructed output series;
they are not a decomposition estimated from data.}
\label{fig:growth-path-diagnosis}
\end{figure}

\Needspace{11\baselineskip}
\begin{redblock}
\textbf{Question: what would your explanation predict next?}
If the slowdown is mainly a closing catch-up gap, restoring yesterday's
investment share need not restore yesterday's growth. If it is mainly weak
utilization, installing more capacity may not solve the immediate problem.
If it reflects a weaker innovation mechanism, a temporary spending recovery
does not by itself repair that mechanism. State a counterfactual and follow
the quantities, decisions, and returns that it changes.
\end{redblock}

\subsection{A longer view: learning to exercise economic judgment}
\label{sec:growth-judgment}

Knowing how to solve each model is the beginning of economic judgment.
The harder task is to decide which relationship deserves attention in an
economy where many things change together. The preceding models help us
practice that judgment: they reveal different costs of moving resources
through time, across people, and toward new uses. Read them together and
a few habits become especially useful.

\subsubsection{Look for the constraint whose relaxation would change behavior}

Suppose the region has both idle factories and small firms that want to
expand. Calling it short of capital or burdened with excess capital is too
quick. Equipment may be in the wrong firms, in the wrong locations, or
suited to yesterday's products. Some producers may lack customers; others
may have orders but lack working funds or skilled workers. An aggregate
capital stock can conceal all of these situations.

Ask what would happen if one particular obstacle were relaxed. Would
cheaper finance lead a firm to produce more, or would it merely refinance
existing debt? Would an additional technician allow equipment to operate,
or would output still be limited by sales? The value of relaxing a constraint
is the intuition behind its shadow price. Observed interest rates, wages,
and profits can help, but taxes, risk, market power, and contractual
restrictions can separate them from social marginal values.

This is why the models should sharpen the next question asked of a firm
or household. ``Investment is low'' names an outcome. A diagnosis identifies
the decision that would change, the obstacle preventing it, and the evidence
that would distinguish that obstacle from another one.

\subsubsection{Recognize that success changes the next problem}

An economy that has built basic equipment may need a different next
investment from one that has scarcely begun. Training can make existing
machines useful; capable engineers can make adaptation feasible; adaptation
can bring production close enough to existing best practice that further
progress requires discovery. These possibilities explain why a successful
policy need not remain the best use of resources indefinitely.

The lesson is not that capital ceases to matter when ideas appear. Research
needs laboratories and people, and new methods may require new equipment.
The relevant complement and the return at the margin can change. Solow's
diminishing returns, the skill model's joint accumulation, and the innovation
model's costly discovery each expose a different part of that process.

A long horizon also changes how we read disappointment. A plant installed
under reasonable expectations can become unprofitable after an unforeseen
change in technology or demand. Its later loss does not establish that its
original decision was irrational. Conversely, repeated losses followed by
new investment require an explanation of incentives, financing, and exit.
Distinguish a bad outcome from a bad decision, then ask whether the decision
rule learns from the outcome.

\subsubsection{Follow resources, financial claims, and expectations together}

Finance gives someone command over resources and specifies who receives
future payments. It does not by itself create the engineers, machines, or
future goods needed to honor those payments. A financing arrangement can
nevertheless be productive if it lets resources reach a use that was
previously blocked. The real allocation and the financial contract must be
examined together.

For example, a guarantee can make funding cheaper for a factory. Ask who
bears the loss if sales disappoint and whether the guarantee changes project
selection. If losses become public obligations, their cost reappears in
someone's future budget. If cheaper funding enables a genuinely productive
project that could not otherwise proceed, that benefit belongs in the
comparison too. Neither the funding rate nor the spending total settles it.

Expectations connect the contract to current choices. A firm pays for
training when it expects to use the resulting skills. A researcher pursues
an idea when future rewards justify its present cost. A household saves
against the consumption opportunities it expects assets to provide.
Institutions matter here through concrete features: the reliability of
contracts, the rules for entry and exit, the treatment of losses, and the
credibility of future policy. Naming ``institutions'' without specifying
such a channel leaves the mechanism unexplained.

\subsubsection{Judge a transition by the people who must live through it}

A better eventual allocation does not tell us how to reach it. Capital
is installed in particular places; skills belong to particular people;
retirement claims were acquired under particular expectations. Moving
toward a new production structure can strand some of those investments.
The aggregate gain and the individual loss can both be real.

Ramsey teaches us to count consumption sacrificed before later benefits
arrive. OLG teaches us that the people who sacrifice and the people who
benefit may be different. An allocation model shows a production gain;
it does not automatically provide the taxes, transfers, mobility, or
information needed to compensate the losers. Even when compensation is
feasible on paper, its financing and delivery require their own account.

The distinction matters for feasibility as well as fairness. A reform can
be reversed if those bearing its early costs have no credible prospect of
sharing its benefits. Yet preserving every existing claim can prevent
resources from moving to better uses. Evaluate the sequence, the temporary
support, and the commitment behind it. Neither immediate liquidation nor
indefinite protection follows from the growth models alone.

\subsubsection{Ask whether the opportunity can be repeated}

Repairing a bridge, moving equipment to a better producer, and inventing
a new method can all raise output. Their contribution to future growth
depends on what remains to be done after the first gain. Once the bridge
works or the equipment is well allocated, repeating the same action need
not yield the same benefit. Research can continue to find new opportunities,
but its resource requirements and difficulty may change.

For every growth proposal, ask what the economy will accumulate or discover
next, what supports the return to that next step, and who will pay for it.
AK makes continued returns explicit as an assumption. Human-capital and
innovation models give those returns a more specific mechanism. The
semi-endogenous experiment shows how a change in research difficulty can
alter the long-run conclusion even while research remains valuable.

This also prevents an excessive preference for growth-rate effects. A
lasting improvement in health, consumption, or productive efficiency can
be worthwhile without making growth permanently faster. The welfare
question concerns the opportunities people gain over their lives, including
their costs. A growth rate is one description of those opportunities.

\subsubsection{Give a conditional judgment and say what would change it}

An experienced economist should be able to state a view clearly while
remaining explicit about what could overturn it. Begin with the leading
mechanism, explain why it fits the available observations, and identify
the strongest competing explanation. Then ask which decision or observation
would separate them. More equations cannot substitute for that comparison.

If orders are weak but productive capacity remains usable, the immediate
question differs from one in which orders are strong and capable producers
cannot expand. If both problems coexist, the policy sequence and its
effects on future capacity must be considered together. An intervention
that improves spending today may also change debt, entry, training, and
the willingness to invest tomorrow.

Some decisions can be revised cheaply; others commit resources for years.
Where learning is possible and delay is not too costly, a limited experiment
with stated evaluation criteria may be preferable to an irreversible
expansion. Where delay destroys productive relationships or skills, waiting
also has a cost. The comparison requires the size of the potential gain,
the consequences of being wrong, and the value of information obtained
along the way.

\Needspace{11\baselineskip}
\begin{greenblock}
\textbf{A useful judgment has a mechanism and a revision rule.}
State which constraint appears to matter, whose behavior would change if
it were relaxed, how the change reaches production and consumption, and who
bears the transition cost. Then state what evidence would make you revise
the diagnosis. The purpose of experience is to make these connections
clearer, including when the available evidence cannot yet choose between
plausible explanations.
\end{greenblock}

\Needspace{20\baselineskip}
\subsection{What we carry into fluctuations}

Growth models leave us with inherited stocks, resource constraints, decision
rules, and a benchmark path. Fluctuation models ask how production and choices
respond when opportunities or spending conditions change over time. The
same physical accounting continues to hold. In a closed, one-good economy
without government, in discrete time,
\[
Y_t=A_tF(K_t,\ell_t),\qquad
K_{t+1}=(1-\delta)K_t+Y_t-C_t,
\]
where $A_t$ is productivity, $\ell_t$ is labor input, and $C_t$ is aggregate
consumption. Here $A_t$ denotes a multiplicative productivity shifter, rather
than necessarily the labor-efficiency index in Solow. An inherited $K_t$
cannot jump because households revise their expectations; consumption,
hours, utilization where modeled, and investment can respond sooner.

The RBC benchmark adds changing real opportunities and intertemporal
responses, with flexible prices and market clearing. New Keynesian models
then add nominal adjustment restrictions and policy transmission. In neither
case is a statistical trend automatically productive capacity, an efficient
allocation, or the model's natural-output path. Those are different objects
whose relationship must be established.

The region may also import equipment, export its products, or borrow abroad.
Our closed-economy resource examples deliberately abstracted from those
transactions. Adding them changes the goods and asset accounts and introduces
relative prices and external financing. The open-economy part will develop
those links; the domestic saving rate then need not equal the domestic
investment rate.

The region's factories, skills, ideas, and age structure remain relevant
when demand falls or productivity changes. Conversely, a recession can
change those stocks and future opportunities. Moving to fluctuations changes
the question and the adjustment mechanisms under examination; it does not
require assuming that all shocks wash out or that long-run growth happens
outside the economy's decisions.

When reading a new episode, first measure the outcome, then identify the
margin of adjustment, the decision-maker, and the missing constraint.
Choose a model that closes those relationships. Its value is the mechanism
it makes explicit and the observations it helps us interpret.
